% 10
\begin{frame}[t]{方向 B：扰动响应中哪些部分可预测？}
\vspace{.1cm}{\color{DarkBlue}\normalsize\bfseries 把观测到的响应拆成扰动间共享的系统性变化和扰动特异信号，分别评估基线与模型；\\[.08cm]
输出一张“哪些扰动、哪些基因模块可预测”的地图，供后续实验选择目标。\par}
\vspace{.12cm}
\centering
\begin{tikzpicture}[nd/.style={box,inner sep=6pt,text width=2.2cm,minimum height=.95cm}]
 \node[nd,fill=PaleBlue] (obs) at (0,0) {\ico{\faWaveSquare}\enspace\textbf{观测响应}\\\scriptsize $\Delta\bar{x}$，来自公开数据};
 \node[nd] (sh) at (3.7,.82) {\textbf{共享的系统性变化}\\\scriptsize 多数扰动都有的方向};
 \node[nd] (sp) at (3.7,-.82) {\textbf{扰动特异信号}\\\scriptsize 区分扰动的部分};
 \node[nd,fill=PaleTeal,draw=DarkBlue,line width=.9pt] (ev) at (7.4,0) {\ico{\faBalanceScale}\enspace\textbf{分层评估}\\\scriptsize 基线 vs 模型\\\scriptsize 按扰动与基因模块};
 \draw[flow] (obs.east) -- (sh.west);\draw[flow] (obs.east) -- (sp.west);
 \draw[flow] (sh.east) -- (ev.west);\draw[flow] (sp.east) -- (ev.west);
 \begin{scope}[shift={(10.45,-.66)}]
  \foreach \i in {0,...,4}{\foreach \j in {0,...,3}{%
    \pgfmathtruncatemacro{\shd}{12+24*mod(\i*2+\j*3,4)}%
    \fill[DarkBlue!\shd] (\i*.32,\j*.32) rectangle ++(.3,.3);}}
  \node[font=\tiny,text=Muted,anchor=north] at (.8,-.06) {基因模块（示意）};
  \node[font=\tiny,text=Muted,anchor=south,rotate=90] at (-.14,.64) {扰动};
  \node[font=\scriptsize\bfseries,anchor=south] at (.8,1.36) {可预测性地图};
  \coordinate (mapw) at (-.32,.64);
 \end{scope}
 \draw[flow] (ev.east) -- (mapw);
 \node[font=\scriptsize,text=Muted,align=center,anchor=north] at (11.25,-1.38) {\faClipboardList\enspace 给合作者的\\可信 / 不可信预测清单};
\end{tikzpicture}
\remark{\ico{\faLightbulb}\enspace 否定结果同样有信息量：若特异信号对所有模型都不可预测，资源应投向数据而非模型。}
\source{分解与扰动区分评估参考 [11] Systema；基线参考 [10]。数据与划分与方向 A 完全共用。候选研究设想，尚未实现。}
\end{frame}
